\documentclass[12pt]{article}
\usepackage{amsmath, amssymb}
\usepackage{geometry}
\usepackage{booktabs}
\usepackage{array}
\usepackage{enumitem}
\usepackage{titlesec}
\usepackage[dvipsnames, table]{xcolor}
\usepackage{fancyhdr}
\usepackage{tcolorbox}
\usepackage{microtype}

\tcbuselibrary{skins}

\geometry{margin=1in, top=1in, bottom=1in}

% ── Colour palette ───────────────────────────────────────────────────────────
\definecolor{primary}  {RGB}{38,  50,  56}
\definecolor{accent}   {RGB}{0,  188, 212}
\definecolor{lightbg}  {RGB}{224,247,250}
\definecolor{gold}     {RGB}{255,152,  0}
\definecolor{goldbg}   {RGB}{255,243,224}
\definecolor{softgray} {RGB}{245,246,247}
\definecolor{darktext} {RGB}{33,  33,  33}
\definecolor{green1}   {RGB}{0,  150, 136}
\definecolor{greenbg}  {RGB}{224,242,241}

% ── Headers / Footers ────────────────────────────────────────────────────────
\pagestyle{fancy}
\fancyhf{}
\fancyhead[L]{\small\color{primary}\textbf{Applications of Laplace Transform}\;\textcolor{accent}{|}\;Assignment}
\fancyhead[R]{\small\color{darktext}BS Mathematics}
\renewcommand{\headrulewidth}{0pt}
\fancyhead[C]{\color{accent}\rule[-1ex]{\textwidth}{1.2pt}}
\fancyfoot[C]{\small\color{darktext}\thepage}

% ── Section styling ──────────────────────────────────────────────────────────
\titleformat{\section}
  {\large\bfseries\color{primary}}{}
  {0em}{}
  [\color{accent}\rule{\linewidth}{2pt}]
\titlespacing*{\section}{0pt}{20pt}{10pt}

\titleformat{\subsection}
  {\normalsize\bfseries\color{accent}}{}
  {0em}{}

% ── tcolorbox styles ─────────────────────────────────────────────────────────
\newtcolorbox{formulabox}{
  enhanced,
  colback=goldbg, colframe=gold, arc=5pt, boxrule=1.5pt,
  left=12pt, right=12pt, top=10pt, bottom=10pt,
  title={\small\bfseries Key Result / Formula},
  colbacktitle=gold!30, coltitle=darktext,
  attach boxed title to top left={xshift=12pt, yshift=-\tcboxedtitleheight/2},
  boxed title style={arc=3pt, colframe=gold, colback=gold!25, boxrule=1pt},
}

\newtcolorbox{conceptbox}{
  enhanced,
  borderline west={4pt}{0pt}{accent},
  colback=lightbg!55, colframe=white,
  arc=0pt, boxrule=0pt,
  left=12pt, right=8pt, top=6pt, bottom=6pt,
}

\newtcolorbox{examplebox}{
  enhanced,
  colback=softgray, colframe=primary!50, arc=5pt, boxrule=1pt,
  left=12pt, right=12pt, top=8pt, bottom=8pt,
  title={\small\bfseries Worked Example},
  colbacktitle=primary, coltitle=white,
}

\newtcolorbox{answerbox}{
  enhanced,
  colback=greenbg, colframe=green1, arc=5pt, boxrule=1.5pt,
  left=12pt, right=12pt, top=8pt, bottom=8pt,
  title={\small\bfseries Result},
  colbacktitle=green1!80, coltitle=white,
  attach boxed title to top left={xshift=12pt, yshift=-\tcboxedtitleheight/2},
  boxed title style={arc=3pt, colframe=green1, colback=green1!70, boxrule=1pt},
}

\newcolumntype{C}[1]{>{\centering\arraybackslash}p{#1}}
\newcommand{\hcell}[1]{\textbf{\color{white}#1}}

% ── Custom math commands ─────────────────────────────────────────────────────
\newcommand{\Lap}[1]{\mathcal{L}\!\left\{#1\right\}}
\newcommand{\iLap}[1]{\mathcal{L}^{-1}\!\left\{#1\right\}}

\begin{document}

% ════════════════════════════════════════════════════════════════════════════
%  TITLE BANNER
% ════════════════════════════════════════════════════════════════════════════
\thispagestyle{empty}

\begin{tcolorbox}[
  enhanced, arc=0pt, boxrule=0pt,
  colback=primary, colframe=primary,
  left=16pt, right=16pt, top=18pt, bottom=14pt,
  borderline south={4pt}{0pt}{accent},
]
  \centering
  {\huge\bfseries\color{white} Applications of Laplace Transform}\\[4pt]
  {\large\color{accent} Assignment --- Transform Theories}\\[6pt]
  {\normalsize\color{lightbg}
    Shah Abdul Latif University, Khairpur \quad\textbullet\quad Department of Mathematics}
\end{tcolorbox}

\vspace{4pt}

\begin{tcolorbox}[
  enhanced, arc=0pt, boxrule=0pt,
  colback=white, colframe=white,
  borderline north={1.5pt}{0pt}{accent!60},
  borderline south={1.5pt}{0pt}{accent!60},
  left=16pt, right=16pt, top=12pt, bottom=12pt,
]
  \begin{tabular}{p{3.8cm} p{9cm}}
    \textcolor{primary}{\textbf{Student:}}
      & \textbf{Nadir Ali Khan} \quad|\quad Roll No: \textbf{MT-0123-052} \\[4pt]
    \textcolor{primary}{\textbf{Subject:}}
      & \textbf{Transform Theories} \\[4pt]
    \textcolor{primary}{\textbf{Submitted To:}}
      & \textbf{Dr.\ Israr Ahmed Memon} \\[4pt]
    \textcolor{primary}{\textbf{Department:}}
      & \textbf{Department of Mathematics} \\
  \end{tabular}
\end{tcolorbox}

\vspace{10pt}

% ── Table of Contents ────────────────────────────────────────────────────────
\renewcommand{\contentsname}{\color{primary}Table of Contents}
\tableofcontents

\newpage

% ════════════════════════════════════════════════════════════════════════════
\section{Introduction to Laplace Transform}
% ════════════════════════════════════════════════════════════════════════════

\begin{conceptbox}
The \textbf{Laplace Transform} is one of the most powerful mathematical tools used in engineering,
physics, and applied mathematics. Named after \textbf{Pierre-Simon Laplace} (1749--1827), it converts
functions of time into functions of a complex variable $s$.
\end{conceptbox}

\medskip
This transformation simplifies the analysis of linear time-invariant systems by turning differential
equations into algebraic equations. Its applications span electrical engineering, mechanical
engineering, control theory, signal processing, and mathematical physics.

The Laplace Transform acts as a bridge between the \emph{time domain} and the
\emph{frequency (complex) domain}, enabling engineers and mathematicians to solve problems that
would be extremely difficult using classical methods alone.

% ════════════════════════════════════════════════════════════════════════════
\section{Definition and Basic Properties}
% ════════════════════════════════════════════════════════════════════════════

\subsection*{Definition}

The Laplace Transform of $f(t)$, defined for $t \geq 0$, is:

\begin{formulabox}
\[
  \Lap{f(t)} = F(s) = \int_{0}^{\infty} e^{-st}\, f(t)\, dt
\]
where $s = \sigma + j\omega$ is a complex number.
\end{formulabox}

\subsection*{Key Properties}

\begin{conceptbox}
\begin{tabular}{lp{8cm}}
\textbf{Linearity:}         & $\Lap{af(t)+bg(t)} = aF(s)+bG(s)$ \\[4pt]
\textbf{Differentiation:}   & $\Lap{f'(t)} = sF(s) - f(0)$ \\[4pt]
\textbf{Integration:}       & $\displaystyle\Lap{\int_0^t f(\tau)\,d\tau} = \dfrac{F(s)}{s}$ \\[8pt]
\textbf{Time Shifting:}     & $\Lap{f(t-a)\,u(t-a)} = e^{-as}F(s)$ \\[4pt]
\textbf{Frequency Shifting:}& $\Lap{e^{at}f(t)} = F(s-a)$ \\[4pt]
\textbf{Convolution:}       & $\Lap{(f*g)(t)} = F(s)\cdot G(s)$ \\[4pt]
\textbf{Initial Value:}     & $\displaystyle\lim_{t\to 0^+} f(t) = \lim_{s\to\infty} sF(s)$ \\[8pt]
\textbf{Final Value:}       & $\displaystyle\lim_{t\to\infty} f(t) = \lim_{s\to 0} sF(s)$ \\
\end{tabular}
\end{conceptbox}

\subsection*{Common Laplace Pairs}

\begin{center}
\begin{tabular}{C{3.5cm} C{5cm} C{3cm}}
\toprule
\rowcolor{primary}
\hcell{$f(t)$} & \hcell{$F(s) = \Lap{f(t)}$} & \hcell{Region of Conv.} \\
\midrule
\rowcolor{lightbg!50}
$1$            & $\dfrac{1}{s}$                     & $s > 0$ \\[6pt]
$t$            & $\dfrac{1}{s^2}$                   & $s > 0$ \\[6pt]
\rowcolor{lightbg!50}
$e^{at}$       & $\dfrac{1}{s-a}$                   & $s > a$ \\[6pt]
$\sin(\omega t)$ & $\dfrac{\omega}{s^2+\omega^2}$   & $s > 0$ \\[6pt]
\rowcolor{lightbg!50}
$\cos(\omega t)$ & $\dfrac{s}{s^2+\omega^2}$        & $s > 0$ \\[6pt]
$\delta(t)$    & $1$                                & all $s$ \\
\bottomrule
\end{tabular}
\end{center}

\newpage

% ════════════════════════════════════════════════════════════════════════════
\section{Application in Solving Ordinary Differential Equations}
% ════════════════════════════════════════════════════════════════════════════

One of the most classical applications of the Laplace Transform is solving linear ODEs with constant
coefficients, particularly initial value problems (IVPs). The process transforms a differential
equation into an algebraic equation, solves it, then inverts the result.

\subsection*{General Procedure}

\begin{enumerate}[noitemsep]
  \item Take the Laplace Transform of both sides of the ODE.
  \item Apply initial conditions to simplify.
  \item Solve the resulting algebraic equation for $F(s)$.
  \item Apply the Inverse Laplace Transform to find $f(t)$.
\end{enumerate}

\subsection*{Example 1 --- First Order ODE}

\begin{examplebox}
Solve $y' + 3y = 6$, \quad $y(0) = 2$.

\medskip
\textbf{Step 1 --- Apply Laplace Transform:}
\[
  \bigl[sY(s) - y(0)\bigr] + 3Y(s) = \frac{6}{s}
\]
\textbf{Step 2 --- Substitute $y(0) = 2$:}
\[
  Y(s)(s+3) = \frac{6}{s} + 2 = \frac{6+2s}{s}
\]
\textbf{Step 3 --- Partial Fractions:}
\[
  Y(s) = \frac{6+2s}{s(s+3)} = \frac{2}{s} + \frac{0}{s+3}
\]
\textbf{Step 4 --- Inverse Transform:}
\[
  \boxed{y(t) = 2}
\]
\end{examplebox}

\subsection*{Example 2 --- Second Order ODE}

\begin{examplebox}
Solve $y'' + 5y' + 6y = 0$, \quad $y(0)=2,\; y'(0)=-1$.

\medskip
Taking the Laplace Transform:
\[
  \bigl[s^2Y - 2s + 1\bigr] + 5\bigl[sY - 2\bigr] + 6Y = 0
\]
\[
  Y(s^2 + 5s + 6) = 2s + 9 \implies Y(s) = \frac{2s+9}{(s+2)(s+3)}
\]

Partial fractions: $\dfrac{2s+9}{(s+2)(s+3)} = \dfrac{A}{s+2} + \dfrac{B}{s+3}$

\[
  A = \frac{2(-2)+9}{-2+3} = 5, \qquad B = \frac{2(-3)+9}{-3+2} = 3
\]

\[
  \boxed{y(t) = 5e^{-2t} + 3e^{-3t}}
\]
\end{examplebox}

\subsection*{Example 3 --- ODE with Heaviside Step Function}

\begin{examplebox}
Solve $y' + y = u(t-2)$, \quad $y(0) = 0$.

\medskip
Since $\Lap{u(t-2)} = \dfrac{e^{-2s}}{s}$:
\[
  sY + Y = \frac{e^{-2s}}{s} \implies Y(s) = \frac{e^{-2s}}{s(s+1)} = e^{-2s}\left(\frac{1}{s} - \frac{1}{s+1}\right)
\]
\[
  \boxed{y(t) = \bigl[1 - e^{-(t-2)}\bigr]\,u(t-2)}
\]
\end{examplebox}

\newpage

% ════════════════════════════════════════════════════════════════════════════
\section{Application in Electrical Circuit Analysis}
% ════════════════════════════════════════════════════════════════════════════

The Laplace Transform is one of the most natural tools for circuit analysis. It provides a
systematic method to analyze RLC circuits, especially those with initial conditions (stored energy).

\subsection*{Transformed Impedances in the $s$-Domain}

\begin{conceptbox}
\begin{tabular}{lp{9cm}}
\textbf{Resistor:}   & $Z_R = R$ \\[4pt]
\textbf{Inductor:}   & $Z_L = sL$ \quad (initial current $i(0)$ adds source $sLi(0)$) \\[4pt]
\textbf{Capacitor:}  & $Z_C = \dfrac{1}{sC}$ \quad (initial voltage $v(0)$ adds source $v(0)/s$) \\
\end{tabular}
\end{conceptbox}

\subsection*{Example --- Series RLC Circuit}

\begin{examplebox}
$R = 2\,\Omega$, $L = 1\,\text{H}$, $C = 0.5\,\text{F}$, driven by $v(t) = 10u(t)$, zero initial conditions.

\medskip
KVL gives: $\quad i''(t) + 2i'(t) + 2i(t) = 10\delta(t)$

Taking Laplace Transform:
\[
  s^2 I + 2sI + 2I = 10 \implies I(s) = \frac{10}{s^2+2s+2} = \frac{10}{(s+1)^2+1}
\]
Inverse Transform:
\[
  \boxed{i(t) = 10\,e^{-t}\sin(t)\;\text{A}, \quad t \geq 0}
\]
This shows a damped sinusoidal response characteristic of an under-damped RLC circuit.
\end{examplebox}

\subsection*{Transfer Function}

The transfer function $H(s) = V_\text{out}(s)/V_\text{in}(s)$ for an RC low-pass filter:
\[
  H(s) = \frac{1/sC}{R + 1/sC} = \frac{1}{1 + sRC} = \frac{1}{1+s\tau}
\]
where $\tau = RC$ is the time constant. The cut-off frequency is $\omega_c = 1/\tau$.

% ════════════════════════════════════════════════════════════════════════════
\section{Application in Control Systems}
% ════════════════════════════════════════════════════════════════════════════

Control systems engineering relies heavily on the Laplace Transform for analyzing and designing
feedback systems in the $s$-domain.

\subsection*{Closed-Loop Transfer Function}

For a feedback system with plant $G(s)$ and controller $C(s)$:
\[
  T(s) = \frac{C(s)G(s)}{1 + C(s)G(s)}
\]

\subsection*{Stability via Poles}

\begin{conceptbox}
\begin{itemize}[noitemsep]
  \item All poles in the \textbf{left half} of the $s$-plane $\Rightarrow$ \textbf{Stable}
  \item Any pole in the \textbf{right half} $\Rightarrow$ \textbf{Unstable}
  \item Poles on the \textbf{imaginary axis} $\Rightarrow$ \textbf{Marginally stable}
\end{itemize}
\end{conceptbox}

\subsection*{PID Controller}

A PID controller has transfer function:
\[
  C(s) = K_p + \frac{K_i}{s} + K_d s = \frac{K_d s^2 + K_p s + K_i}{s}
\]

\newpage

% ════════════════════════════════════════════════════════════════════════════
\section{Application in Mechanical Vibrations}
% ════════════════════════════════════════════════════════════════════════════

The Laplace Transform is extensively applied in analysing vibrating systems --- spring-mass systems,
beams, and rotating machinery.

\subsection*{Spring-Mass-Damper System}

Equation of motion:
\[
  m\ddot{x}(t) + c\dot{x}(t) + kx(t) = f(t)
\]
Taking Laplace Transform (zero initial conditions):
\[
  (ms^2 + cs + k)X(s) = F(s) \implies H(s) = \frac{X(s)}{F(s)} = \frac{1}{ms^2+cs+k}
\]

\begin{examplebox}
$m = 1\,\text{kg}$, $c = 2\,\text{N\,s/m}$, $k = 5\,\text{N/m}$, $f(t) = 10u(t)$.

\[
  X(s) = \frac{10}{s(s^2+2s+5)} = \frac{2}{s} - \frac{2(s+1)}{(s+1)^2+4} - \frac{2}{(s+1)^2+4}
\]
\[
  \boxed{x(t) = 2 - 2e^{-t}\cos(2t) - e^{-t}\sin(2t)\;\text{ m}}
\]
Steady-state displacement: $x_{ss} = F/k = 10/5 = 2\,\text{m}$ \checkmark
\end{examplebox}

\subsection*{Damping Classification}

\begin{conceptbox}
Natural frequency: $\omega_n = \sqrt{k/m}$ \quad|\quad Damping ratio: $\zeta = c\,/\,(2\sqrt{mk})$
\begin{itemize}[noitemsep, topsep=4pt]
  \item $\zeta < 1$: \textbf{Under-damped} (oscillatory response)
  \item $\zeta = 1$: \textbf{Critically damped} (fastest non-oscillatory)
  \item $\zeta > 1$: \textbf{Over-damped} (slow exponential return)
\end{itemize}
\end{conceptbox}

% ════════════════════════════════════════════════════════════════════════════
\section{Application in Signal Processing}
% ════════════════════════════════════════════════════════════════════════════

The Laplace Transform provides the mathematical foundation for analysing and designing filters,
modulators, and communication systems. It handles transient and exponentially growing signals
that the Fourier Transform cannot.

\subsection*{Convolution Theorem}

For an LTI system with impulse response $h(t)$ and input $x(t)$:
\[
  y(t) = x(t)*h(t) = \int_0^t x(\tau)\,h(t-\tau)\,d\tau
\]

\begin{formulabox}
In the $s$-domain: \quad $Y(s) = X(s)\cdot H(s)$

Convolution in time $\equiv$ multiplication in the $s$-domain.
\end{formulabox}

\subsection*{Filter Design}

A Butterworth low-pass filter of order $n$ with cut-off $\omega_c$:
\[
  |H(j\omega)|^2 = \frac{1}{1 + (\omega/\omega_c)^{2n}}
\]
Poles lie on a circle of radius $\omega_c$ in the left half $s$-plane, ensuring stability and
maximally flat frequency response.

\newpage

% ════════════════════════════════════════════════════════════════════════════
\section{Application in Heat Conduction (PDEs)}
% ════════════════════════════════════════════════════════════════════════════

The Laplace Transform reduces PDEs to ODEs by transforming with respect to time.

\subsection*{Heat Equation}

The one-dimensional heat equation:
\[
  \frac{\partial u}{\partial t} = \alpha\,\frac{\partial^2 u}{\partial x^2}
\]

Taking $\Lap{\cdot}$ with respect to $t$ (let $U(x,s) = \Lap{u(x,t)}$):
\[
  sU(x,s) - u(x,0) = \alpha\,\frac{d^2U}{dx^2}
  \qquad\Longrightarrow\qquad
  \frac{d^2U}{dx^2} - \frac{s}{\alpha}U = -\frac{u(x,0)}{\alpha}
\]

This is now an ODE in $x$, solved by standard methods, then inverted.

\subsection*{Example --- Semi-Infinite Rod}

\begin{examplebox}
Semi-infinite rod ($x \geq 0$), initially at $u=0$, surface held at $u(0,t)=T_0$.

\[
  U(x,s) = \frac{T_0}{s}\,e^{-x\sqrt{s/\alpha}}
\]
\[
  \boxed{u(x,t) = T_0\,\operatorname{erfc}\!\left(\frac{x}{2\sqrt{\alpha t}}\right)}
\]
where $\operatorname{erfc}$ is the complementary error function. This describes heat
penetration from the surface into the rod over time.
\end{examplebox}

\subsection*{Other PDE Applications}

\begin{itemize}[noitemsep]
  \item \textbf{Wave equation:} Modelling vibrations in strings and membranes.
  \item \textbf{Diffusion equation:} Mass transfer in chemical processes.
  \item \textbf{Telegraph equation:} Signal transmission in cables.
  \item \textbf{Beam equations:} Deflection under dynamic loads in structural engineering.
\end{itemize}

% ════════════════════════════════════════════════════════════════════════════
\section{Conclusion}
% ════════════════════════════════════════════════════════════════════════════

The Laplace Transform is far more than a mathematical curiosity --- it is an essential engineering
and scientific tool. From transforming differential equations into algebra, to providing a unified
framework for stability analysis, frequency response, mechanical damping, heat diffusion, and signal
filtering, the Laplace Transform remains indispensable.

Its elegance lies in the duality it creates: complex time-dependent behaviour becomes a static
algebraic object in the $s$-domain. This perspective --- viewing a system's entire history
compressed into a rational function of $s$ --- is one of the most profound insights of applied
mathematics.

\begin{answerbox}
\textbf{Key Takeaways:}
\begin{itemize}[noitemsep, topsep=4pt]
  \item Converts ODEs/PDEs $\to$ algebraic equations in the $s$-domain.
  \item Handles initial conditions automatically via differentiation property.
  \item Convolution in time $\equiv$ multiplication in $s$-domain.
  \item Poles of $H(s)$ determine system stability completely.
  \item Unified framework across circuits, control, mechanics, heat, and signals.
\end{itemize}
\end{answerbox}

\vspace{14pt}

% ════════════════════════════════════════════════════════════════════════════
\section*{References}
% ════════════════════════════════════════════════════════════════════════════

\begin{enumerate}[noitemsep]
  \item Kreyszig, E.\ (2011). \textit{Advanced Engineering Mathematics} (10th ed.). John Wiley \& Sons.
  \item Stroud, K.\ A., \& Booth, D.\ J.\ (2011). \textit{Engineering Mathematics} (7th ed.). Palgrave Macmillan.
  \item Boyce, W.\ E., \& DiPrima, R.\ C.\ (2017). \textit{Elementary Differential Equations and Boundary Value Problems} (11th ed.). Wiley.
  \item Ogata, K.\ (2010). \textit{Modern Control Engineering} (5th ed.). Prentice Hall.
  \item Hayt, W.\ H., Kemmerly, J.\ E., \& Durbin, S.\ M.\ (2012). \textit{Engineering Circuit Analysis} (8th ed.). McGraw-Hill.
  \item Oppenheim, A.\ V., \& Willsky, A.\ S.\ (1997). \textit{Signals and Systems} (2nd ed.). Prentice Hall.
  \item Cengel, Y.\ A.\ (2014). \textit{Heat and Mass Transfer} (5th ed.). McGraw-Hill.
\end{enumerate}

\end{document}
